\documentclass[12pt,a4paper]{article}

% --- PAQUETES ---
\usepackage[utf8]{inputenc}
\usepackage[spanish]{babel}
\usepackage{amsmath}
\usepackage{geometry}
\usepackage[american]{circuitikz}
\usetikzlibrary{babel} % Solución para conflicto circuitikz y babel
\usepackage{float}
\usepackage{booktabs}
\usepackage{siunitx}
\usepackage{fancyhdr}
\usepackage[colorlinks=true, linkcolor=blue, urlcolor=blue]{hyperref}

% --- CONFIGURACIÓN ---
\geometry{a4paper, margin=1in}
\sisetup{output-decimal-marker={,}}
\renewcommand{\figurename}{Figura}

% --- ENCABEZADO Y PIE DE PÁGINA ---
\pagestyle{fancy}
\fancyhf{}
\fancyhead[R]{Evaluación de Circuitos DC}
\fancyfoot[C]{\thepage}
\renewcommand{\headrulewidth}{0.4pt}

\title{Evaluación de Examen 1: Fundamentos de Circuitos DC}
\author{Estudiante: Sebastián Pulido \\ \small{C.I. 89.552.210} \\ \vspace{0.5cm} \large{Revisado por: Prof. César Rodríguez}}
\date{20 de mayo de 2026}

\begin{document}

\maketitle

\section*{Resumen de Calificación}

\begin{center}
    \huge \textbf{Calificación Final: 2,50 / 10,00}
\end{center}

\vspace{1cm}

\begin{table}[H]
    \centering
    \caption{Desglose de Puntuación por Problema}
    \label{tab:calificaciones}
    \begin{tabular}{@{}lcc@{}}
    \toprule
    \textbf{Problema} & \textbf{Puntuación Máxima} & \textbf{Puntuación Obtenida} \\
    \midrule
    1. Asociación de Resistencias & 3,33 & 1,00 \\
    2. Análisis Nodal & 3,33 & 0,00 \\
    3. Análisis de Mallas & 3,33 & 1,50 \\
    \midrule
    \textbf{TOTAL} & \textbf{10,00} & \textbf{2,50} \\
    \bottomrule
    \end{tabular}
\end{table}

\newpage

\section{Evaluación Detallada}

\subsection{Problema 1: Asociación de Resistencias}
\addcontentsline{toc}{subsection}{Problema 1: Asociación de Resistencias}

\begin{itemize}
    \item \textbf{Calificación: 1,00 / 3,33}
    \item \textbf{Comentario:} El estudiante comete un error fundamental al interpretar la topología del circuito. Aunque los dos primeros pasos de cálculo son correctos, el resto del procedimiento se basa en una conexión incorrecta de los componentes, llevando a un resultado erróneo.
\end{itemize}

\paragraph{Análisis del Procedimiento:}
\begin{enumerate}
    \item \textbf{Parcialmente Correcto:} El estudiante calcula correctamente el paralelo de $R_2$ y $R_3$ ($R_{23} = 10\Omega$) y su suma en serie con $R_1$ ($R_{123} = 15\Omega$).
    \item \textbf{Incorrecto:} A partir de este punto, el estudiante asume erróneamente que $R_4$ y $R_5$ están en paralelo, y que este bloque resultante está en serie con el bloque anterior. La topología correcta del circuito es $((R_1 + (R_2 \parallel R_3)) \parallel R_4) + R_5$.
    \item \textbf{Resultado Incorrecto:} El resultado final del estudiante es $\mathbf{R_{eq} = 17,14\Omega}$. El resultado correcto es $\mathbf{10\Omega}$.
\end{enumerate}

\begin{figure}[H]
    \centering
    \begin{circuitikz}[american]
        % Define nodes for the main parallel block
        \coordinate (P_start) at (0,2);
        \coordinate (P_end) at (7,2);

        % Upper branch: R1 + (R2||R3)
        \draw (P_start) to[R, l=$R_1$ (5$\Omega$)] (2,2);
        \draw (2,2) to[short, *-] ++(0,1) to[R, l=$R_2$ (20$\Omega$)] ++(4,0) to[short, -*] (6,2);
        \draw (2,2) to[short, *-] ++(0,-1) to[R, l_=$R_3$ (20$\Omega$)] ++(4,0) to[short, -*] (6,2);
        \draw (6,2) -- (P_end);

        % Lower branch: R4
        \draw (P_start) -- (0,0) to[R, l_=$R_4$ (15$\Omega$)] (7,0) -- (P_end);

        % R5 in series at the end
        \draw (P_end) to[R, l={$R_5$ (2,5$\Omega$)}] (9,2);

        % Terminals
        \draw (P_start) node[left] {A};
        \draw (9,2) node[right] {B};
    \end{circuitikz}
    \caption{Circuito del Problema 1.}
\end{figure}

\newpage

\subsection{Problema 2: Análisis Nodal}
\addcontentsline{toc}{subsection}{Problema 2: Análisis Nodal}

\begin{itemize}
    \item \textbf{Calificación: 0,00 / 3,33}
    \item \textbf{Comentario:} El estudiante no aplica el método de análisis nodal. No se formula ninguna ecuación de nodo y los cálculos presentados son conceptualmente incorrectos y no conducen a la solución del problema.
\end{itemize}

\paragraph{Análisis del Procedimiento:}
El estudiante no plantea las ecuaciones de la Ley de Corrientes de Kirchhoff (LKC). El único cálculo visible, $I = 5A / 2\Omega = 2.5A$, es una aplicación incorrecta de la Ley de Ohm que no tiene relación con el método solicitado.

\paragraph{Solución Correcta:}
Las ecuaciones correctas son:
\begin{itemize}
    \item \textbf{Nodo A:} $5 = \frac{v_A}{2} + \frac{v_A - v_B}{2} \implies \mathbf{2v_A - v_B = 10}$
    \item \textbf{Nodo B:} $\frac{v_A - v_B}{2} = \frac{v_B}{4} \implies \mathbf{2v_A - 3v_B = 0}$
\end{itemize}
La resolución de este sistema da como resultado $\mathbf{v_A = 7,5\,V}$ y $\mathbf{v_B = 5\,V}$.

\begin{figure}[H]
    \centering
    \begin{circuitikz}[american]
        \draw (0,0) node[ground]{} to[short] (0,2)
            to[isource, l=5A] (0,4)
            to[short] (2,4) node[above, xshift=-0.8cm]{Nodo A ($v_A$)};
        \draw (2,4) to[R, l=2$\Omega$] (2,0);
        \draw (2,4) to[R, l_=2$\Omega$] (4,4) node[above, xshift=0.5cm]{Nodo B ($v_B$)};
        \draw (4,4) to[R, l=4$\Omega$] (4,0);
        \draw (0,0) -- (4,0);
    \end{circuitikz}
    \caption{Circuito del Problema 2.}
\end{figure}

\newpage

\subsection{Problema 3: Análisis de Mallas}
\addcontentsline{toc}{subsection}{Problema 3: Análisis de Mallas}

\begin{itemize}
    \item \textbf{Calificación: 1,50 / 3,33}
    \item \textbf{Comentario:} El estudiante logra plantear correctamente el sistema de ecuaciones en un segundo intento, pero la resolución algebraica es extremadamente desordenada y contiene múltiples errores de signo. Milagrosamente, llega a la relación correcta para $I_2$, pero no completa el problema al no calcular la corriente solicitada en el resistor central.
\end{itemize}

\paragraph{Análisis del Procedimiento:}
\begin{itemize}
    \item \textbf{Planteamiento de Ecuaciones:} Tras un primer intento con errores, el estudiante plantea correctamente el sistema de ecuaciones:
    \begin{itemize}
        \item \textbf{Malla 1:} $9I_1 - 5I_2 = 12$
        \item \textbf{Malla 2:} $5I_1 - 7I_2 = 6$ (o su equivalente $-5I_1 + 7I_2 = -6$)
    \end{itemize}
    \item \textbf{Resolución:} El método de eliminación está plagado de errores de signo. Sin embargo, de forma inexplicable, el estudiante obtiene la ecuación correcta $38I_2 = 6$.
    \item \textbf{Resultado Incompleto:} Se calcula correctamente $I_2 \approx 0,16\,A$, pero no se calcula $I_1$ ni la corriente final solicitada, que es $I_{5\Omega} = I_1 - I_2$.
\end{itemize}

\begin{figure}[H]
    \centering
    \begin{circuitikz}[american]
        \draw (0,0) node[ground]{}
            to[V, l=12V] (0,3)
            to[R, l=4$\Omega$] (3,3)
            to[R, l=5$\Omega$, i>=$I_{5\Omega}$] (3,0) -- (0,0);
        \draw (3,3) to[R, l=2$\Omega$] (6,3)
            to[V, l=6V, invert] (6,0) -- (3,0);
        \draw[->, thick, red] (1.5, 0.7) arc (-90:270:0.8) node[right, pos=0.25]{$I_1$};
        \draw[->, thick, red] (4.5, 0.7) arc (-90:270:0.8) node[right, pos=0.25]{$I_2$};
    \end{circuitikz}
    \caption{Circuito del Problema 3.}
\end{figure}

\newpage

\section*{Comentarios Finales}

El estudiante Sebastián Pulido muestra tener serias dificultades con los conceptos fundamentales de análisis de circuitos y con la ejecución matemática requerida.

\paragraph{Fortalezas:}
\begin{itemize}
    \item Capacidad para realizar cálculos de asociaciones simples (paralelo de $R_2, R_3$).
    \item Logró, eventualmente, plantear las ecuaciones de malla correctas para el problema 3.
\end{itemize}

\paragraph{Áreas Críticas de Mejora:}
\begin{itemize}
    \item \textbf{Interpretación de Circuitos:} Es fundamental aprender a leer correctamente un diagrama esquemático. El error en el problema 1 es grave y demuestra una falta de comprensión de las topologías básicas.
    \item \textbf{Aplicación de Métodos de Análisis:} El estudiante no demostró saber aplicar el análisis nodal (problema 2). Es crucial estudiar y practicar la aplicación sistemática de los métodos.
    \item \textbf{Rigor Algebraico:} La resolución del problema 3 fue caótica y llena de errores. La habilidad para manipular sistemas de ecuaciones de forma limpia y sin errores es indispensable.
    \item \textbf{Atención a la Pregunta:} Es importante leer y responder completamente lo que se solicita en el enunciado (ej. calcular la corriente final en el problema 3).
\end{itemize}

Se recomienda una revisión profunda de los temas de la unidad, comenzando por la correcta identificación de conexiones en serie y paralelo, y practicando intensivamente la formulación y resolución de ecuaciones para análisis nodal y de mallas.

\end{document}